\section{Reproducibility}
\label{sec:reproducibility}

\paragraph{Code and data availability.}
All code, trained model weights, synthgen corpus, and fixture data are
available at:
\begin{center}
  \url{https://github.com/edenduthie/infon}
\end{center}
The exact version used to produce all results in this paper is tagged
\texttt{v0.4.0-paper}.

\paragraph{Installation.}
The study dependencies install into a standard Python~3.11+ environment
via:
\begin{verbatim}
pip install -e ".[study]"
\end{verbatim}
No GPU is required.  SPLADE-tiny is downloaded automatically from
HuggingFace on first use (17~MB).

\paragraph{Reproducing all results.}
A single shell script reproduces the full evaluation matrix:
\begin{verbatim}
bash paper/reproducer.sh
\end{verbatim}
This script: (1) downloads and verifies the HoVer, AVeriTeC, and SciFact
datasets; (2) runs all six systems over the evaluation splits; (3) writes
fixture JSON files to \texttt{paper/tests/fixtures/}; and (4) generates
all figures as PDF files in \texttt{paper/figures/}.

\paragraph{Expected wall-clock time.}
The full evaluation matrix requires 4--8 hours on a modern CPU (2024
vintage, 8+ cores) using the default MCTS configuration.  Figure
generation only (skipping system evaluation and loading pre-computed
fixture data) completes in approximately 10 minutes.

\paragraph{Compute resources.}
All experiments were run on CPU only; no GPU was used at any stage.  Peak
memory consumption is approximately 4~GB RAM.  Dataset caches require
approximately 2~GB of disk storage.

\paragraph{Verification.}
After running \texttt{reproducer.sh}, the numbers audit and figure tests can
be verified with:
\begin{verbatim}
pytest paper/tests/ -v
make lint
\end{verbatim}
All tests must pass with exit code~0.
